\section[Bibliothèque numérique]{Histoire de la bibliothèque numérique}
L'élaboration de la bibliothèque numérique CujasNum ne s'est pas faite de manière linéaire, mais résulte d'une succession de jalons éditoriaux et scientifiques qui ont progressivement défini l'identité numérique de l'établissement.

\subsection{Les balbutiements de la numérisation à Cujas (2001-2007)}
Malgré la richesse des collections conservées dans la réserve de la Bibliothèque Cujas, celle-ci ne rencontre pas immédiatement le succès escompté auprès des usagers. Frédérique Joannic-Seta corrobore ce constat, attribuant ce manque d'attrait à \enquote{ses difficultés d'ouverture de consultation qui entraînent une faible utilisation au regard de sa réelle richesse} \footcite{joannic-setaEtudePrealableOperation1999}. Par ailleurs, de nombreux volumes se trouvent alors en mauvais état, ayant subi aussi bien des \enquote{inondations exceptionnelles dans les magasins que des dégradations quotidiennes suscitées par les manipulations et les conditions de conservation, longtemps médiocres} \footcite{joannic-setaEtudePrealableOperation1999}. 

En 1998, soucieuse de valoriser cette réserve patrimoniale, la Bibliothèque Cujas décide d'engager la numérisation de ses collections. Son choix est dicté par \enquote{un souci de conservation des documents par la mise à disposition des lecteurs d'un support de substitution de qualité} \footcite{joannic-setaEtudePrealableOperation1999}. Le second objectif vise à faciliter l'accès aux fonds pour permettre aux lecteurs \enquote{une meilleure appropriation du texte en leur proposant des facilités de reproduction} \footcite{joannic-setaEtudePrealableOperation1999}. 

L'un des premiers chantiers de numérisation de la Bibliothèque Cujas remonte aux années 2000 et 2001. Durant cette phase pionnière, l'établissement fait numériser 17 titres représentant environ 22 volumes, une tâche externalisée auprès de la société Safig \footcite{2006ProgrammeNumerisation}. La bibliothèque commande alors une numérisation \enquote{en mode image, avec une navigation grâce à la pagination}\footnote{\nameref{annexe_A_rapport_audit}.}. Séduite par l'offre du prestataire qui propose de prendre en charge l'hébergement et la diffusion en ligne des documents sur ses propres serveurs, la Bibliothèque Cujas choisit de lui déléguer l'ensemble de cette chaîne de stockage et de mise à disposition. 

Cette stratégie se révèle désastreuse : la société est victime d'une attaque sur ses serveurs, entraînant la perte irrémédiable d'une majeure partie des documents numérisés. Face à ce sinistre, la bibliothèque s'efforce de récupérer ce qu'elle peut \footcite{2006ProgrammeNumerisation}. Le prestataire Safig restitue ce qui lui reste \enquote{sous forme de CD-ROM}. À la réception des CD-ROM, le personnel de Cujas tente de déterminer s'il est possible \enquote{de réutiliser du contenu}. Cette tentative de sauvetage s'avère toutefois vaine: les fichiers réceptionnés se révèlent irrécupérables, voire corrompus pour certains, contraignant l'établissement à prendre la décision de renumériser intégralement le corpus \footnote{\nameref{annexe_A_rapport_audit}.}.

Le choix de la Bibliothèque Cujas est d'orienter d'emblée la valorisation de son contenu vers \enquote{un lectorat universitaire de haut niveau} \footcite{balley4ExempleMediation2012}. Noëlle Balley, alors chargée de mission pour le patrimoine et future responsable du département des Monographies de la Bibliothèque Cujas, et son adjoint de l'époque, Sébastien Dalmon, attribuent cette volonté à la spécialisation extrême du fonds: à leurs yeux, ce contenu \enquote{souvent aride} s'avère peu propice à une \enquote{valorisation orientée vers le grand public} \footcite{balley4ExempleMediation2012}. L'établissement privilégie donc des corpus spécifiquement sélectionnés pour ce public académique, en mettant à sa disposition des \enquote{outils facilitant leur travail sur les textes}\footcite{balley4ExempleMediation2012}. Toutefois, ils soulignent que malgré les moyens mis en oeuvre pour s'adresser à ce lectorat élitiste, il convient de ne pas \enquote{négliger le grand public} \footcite{balley4ExempleMediation2012}.

\subsection{La liste Pfister-Roumy}
Afin de valoriser la richesse de ce patrimoine, la Bibliothèque Cujas conçoit le projet de développer sa propre bibliothèque numérique, une ambition émergeant en même temps que la numérisation des cours de Jean Carbonnier \footcite{2006ProgrammeNumerisation}. À l'origine, l'établissement souhaite s'insérer dans une dynamique nationale en postulant à l'appel à projets de Persée, dans l'espoir de devenir un pôle de réplication \footcite{2006ProgrammeNumerisation}. C'est dans cette perspective que la bibliothèque acquiert un scanner Kirtas destiné à créer sa chaîne de traitement interne. Bien que cette candidature n'ait finalement pas été retenue, elle sera le point de départ de la réflexion de la bibliothèque souhaitant capitaliser sur cette acquisition \footcite{2007RefonteProgramme}. 

La direction de la bibliothèque Cujas change alors de stratégie.. Elle choisit de s'associer étroitement avec deux éminents  historiens du droit de la Faculté de Droit de Paris, les professeurs Laurent Pfister et Franck Roumy \footcite{2006ProgrammeNumerisation}. Cette collaboration inédite entre les professionnels des bibliothèques et les enseignants-chercheurs permet d'élaborer une première liste de référence scientifique, désignée plus tard sous le nom de \enquote{liste Pfister-Roumy} \footcite{2007RefonteProgramme}.

La numérisation de la collection Jean Carbonnier est la première opération de numérisation interne de la Bibliothèque Cujas. Elle a donc permis de tester \enquote{en grandeur réelle, la chaîne de numérisation qu'elle mettait en place pour son propre programme} \footcite{2007RefonteProgramme}. Face au succès de cette première numérisation, les équipes la réemploient au moment de son deuxième plus gros projet de numérisation, la liste \enquote{Pfister-Roumy}.

Contrairement au projet Carbonnier, ce nouveau projet ne se concentre plus sur des cours polycopiés, mais vise à numériser en priorité  les grands textes fondamentaux et les sources doctrinales de l'histoire du droit.

L'objectif de la liste Pfister-Roumy est de sélectionner des monographies et des périodiques issus \enquote{uniquement des collections de Cujas}. Cette sélection porte sur les \enquote{fondamentaux du droit à numériser en priorité} \footcite{2006ProgrammeNumerisation}. À l'origine, les deux professeurs recensent \enquote{471 titres fondamentaux pour la recherche en histoire du droit}, couvrant une vaste période allant du \enquote{XVe siècle au XXe siècle, soit 1462 volumes et 1 001 559 pages}\footcite{2006ProgrammeNumerisation}. À cet ensemble patrimonial, les professeurs adjoignent deux périodiques de référence: \enquote{Duvergier et le Bulletin des arrêts de la Cour de cassation (Chambre criminelle et Chambres civiles) jusqu'en 1950.}\footcite{2006ProgrammeNumerisation}. Ce projet est initialement pensé lors de la candidature de Cujas pour devenir pôle de réplication de Persée, se caractérise par des ambitions techniques élevées, reposant sur une \enquote{numérisation automatique non descriptive, en mode image océrisé, avec une recherche plein texte et la structuration documentaire à granularité fine, emploi de la DTD TEI, réalisée exclusivement en interne} \footcite{2006ProgrammeNumerisation}.

L'abandon de ce projet de réplication oblige la Bibliothèque Cujas à repenser son programme de numérisation, contraignant l'établissement à revoir à la baisse le nombre d'ouvrages sélectionnés. Le corpus initial est ainsi réduit, passant de 471 ouvrages à 436 ouvrages \footcite{2007RefonteProgramme}. En effet, à la même époque, la \ac{BNF} développe en parallèle un projet scientifique similaire, intitulé la liste Kerbrat. Élaborée sous la direction du professeur Yann Kerbrat et d'un groupe de quatorze universitaire, cette sélection a pour objectif \enquote{de créer une liste de références à partir des collections de la BNF} dans le domaine du droit \footcite{2007RefonteProgramme}. Dans un soucis d'optimisation des moyens financiers et de cohérence scientifique, une concertation étroite s'établit entre les deux institutions afin de coordonner leurs campagnes respectives et d'éviter les doublons. Ainsi, lorsque la Bibliothèque Cujas conserve un ouvrage identique à celui de la \ac{BNF}, il lui est néanmoins permis de le numériser, mais à la condition expresse de sélectionner une édition différente de celle de la \ac{BNF}. Certains titres, impossible à substituer par une autre édition, sont alors purement et simplement retirés de la liste de Cujas \footnote{\nameref{annexe_A_rapport_audit}}. 

\subsection{Les cours de droit}
Les 5 et 6 novembre 2008, la Bibliothèque Cujas organise un \enquote{colloque international en collaboration avec le Sénat},portant sur le doyen Jean Carbonnier. Noëlle Balley et Sébastien Dalmon qualifient cette initiative de \enquote{première entreprise de numérisation effective [qui] a concerné, pour des raisons conjoncturelles, les polycopiés du doyen Jean Carbonnier}\footcite{balley4ExempleMediation2012}. La numérisation de ces polycopiés a été réalisée à \enquote{l'occasion du centenaire de la naissance de ce grand juriste civiliste} \footcite{balley4ExempleMediation2012}. À l'origine, ce projet résidait dans une \enquote{participation modeste} de la Bibliothèque Cujas à \enquote{un projet éditorial}, initié par l'anthropologue du droit Raymond Verdier. Son objectif principal était alors \enquote{de rassembler les écrits de Jean Carbonnier introuvables, inédits ou dispersés}\footcite{balley4ExempleMediation2012}. La valorisation de ce projet scientifique revêt diverses formes, comme la tenue de colloques au Sénat, des actions de communication, une exposition virtuelle.

Parallèlement à la tenue des colloques scientifiques, la Bibliothèque Cujas prend l'initiative d'organiser une exposition virtuelle d'envergure. L'établissement fait ainsi le choix \enquote{de mettre en ligne une exposition virtuelle consacrée à Jean Carbonnier, constatant l'absence, sur Internet, de ressources suffisamment fiables et complètes sur le Doyen}\footcite{balley4ExempleMediation2012}. L'objectif est à la fois de faire connaître \enquote{une grande figure du droit contemporain, méconnue du grand public, mais aussi à répondre aux attentes des spécialistes}\footcite{balley4ExempleMediation2012}. 

Pour mener à bien cette valorisation culturelle, la Bibliothèque Cujas s'appuie opportunément sur son expérience acquise en matière d'éditorialisation numérique. L'établissement a en effet déjà réalisé une exposition sur\enquote{le Bicentenaire du Code civil (1804-2004) en collaboration avec la Cour de cassation et l'Ordre des avocats}\footcite{balley4ExempleMediation2012}. Cette antériorité partenariale permet aux équipes de Cujas de mobiliser un savoir-faire technique pour concevoir une exposition virtuelle, entièrement dédiée à la mise en valeur scientifique et mémorielle de Jean Carbonnier \footcite{2007RefonteProgramme}.

L'exposition virtuelle consacrée à Jean Carbonnier offre un \enquote{accès à treize polycopiés de cours du doyen Carbonnier professés entre 1958 et 1974, portant sur la sociologie juridique ou le droit civil} \footcite{balley4ExempleMediation2012}. Ce corpus de documents, qui constitue la première collection intégrée à la bibliothèque numérique, est aujourd'hui hébergé sur la plateforme actuelle -- anciennement désignée sous le nom de CujasNum -- sous l'intitulé \enquote{Cours de droit Carbonnier}\footcite{SciencesJuridiqueslesSciences}. Cette entreprise représente un chantier d'une rare intensité pour l'établissement. En effet, les documents sont numérisés \enquote{en mode texte, avec une relecture intégrale} \footcite{balley4ExempleMediation2012}. Soucieuse d'offrir le \enquote{meilleur accès aux informations contenues dans le texte}, la Bibliothèque Cujas prend le parti d'assurer elle-même la correction complète de l'\ac{OCR}, ses équipes rectifiant manuellement  chaque coquille de reconnaissance afin d'atteindre un taux de précision de 100\% \footnote{Voir l'annexe \ref{annexe_A_rapport_audit}.}. 

Cette numérisation intense amène à la création d'une \enquote{plate-forme de recherche [...] par l'ingénieur informatique de la bibliothèque; des sommaires dynamiques ont également été créés}\footcite{balley4ExempleMediation2012}. Afin de s'adapter au mieux aux exigences des chercheurs, l'ergonomie de l'interface de recherche est profondément améliorée. Elle permet désormais de croiser les \enquote{recherches sur les métadonnées, la table des matières et le plein texte, soit en favorisant la sérendipidité grâce à une navigation par listes d'auteurs ou de titres, ou encore à travers un thésaurus spécialement conçu pour la bibliothèque numérique} \footcite{Arabesques_69_12_dossier_balleypdfEspaceTravail}. Pour parfaire ce dispositif et optimiser le confort de lecture, l'établissement y adjoint plus tard une visionneuse performante, autorisant un accès \enquote{à une page précise sans télécharger tout le document}\footnote{\nameref{annexe_A_rapport_audit}.}.